\documentclass[runningheads]{llncs}
\usepackage[T1]{fontenc}

% graphics packages
\usepackage{graphicx}
\usepackage{svg}
\usepackage{tikz}
\usepackage{standalone}

% math packages
\usepackage{amsmath,amsfonts,amssymb,bbm}
\usepackage{mathtools}
\usepackage{ebproof}
\usepackage{algorithm}
\usepackage{algpseudocode}

% ref packages
\usepackage{hyperref}
\usepackage{cleveref}
\usepackage{listings}
\usepackage[style=authoryear,backend=biber,natbib=true,maxcitenames=1]{biblatex}

% other packages
\usepackage{todonotes}
\usepackage{tabularx}
\usepackage{comment}
\usepackage{stackengine}
\usepackage{csquotes}


% isabelle snippets, see https://astahfrom.github.io/cartouches/

\usepackage{isabelle, isabellesym}
\isabellestyle{it} % optional
\newcommand{\SNIPTEXT}[1]{#1} % remove #1 to ignore text commands
\newcommand{\SNIP}[2]{\expandafter\newcommand\csname snippet--#1\endcsname{#2}}
\input{isa_snippets} % assuming the snippets are saved as snips.tex
\setstackEOL{\\}
\usetikzlibrary{positioning,quotes,calc,shapes.callouts}

\def\isadelimtheory{}\def\endisadelimtheory{}
\def\isatagtheory{}\def\endisatagtheory{}
\def\isadelimML{}\def\endisadelimML{}
\def\isadelimdocument{}\def\endisadelimdocument{}
\def\isatagML{}\def\endisatagML{}
\def\isatagdocument{}\def\endisatagdocument{}
\def\isafoldML{}
\def\isadelimproof{}\def\endisadelimproof{}
\def\isatagproof{}\def\endisatagproof{}
\def\isafoldproof{}

\raggedbottom

% \GetSnip{name} gets the snippet defined by `name' or produces a warning that it does not exist.
\newcommand{\GetSnip}[1]{%
	\ifcsname snippet--#1\endcsname%
	\csname snippet--#1\endcsname%
	\else%
	\PackageWarning{snips}{Snippet ``#1'' is undefined.}%
	\emph{Warning: Snippet ``#1'' is undefined.}%
	\fi%
}

% \RawCartouche{name}{line}{n}: content of cartouche `n' on line `line' of snippet `name'
\newcommand{\RawCartouche}[3]{\GetSnip{#1-#2-#3}}

% \Cartouche{name}{line}{n}: enclosed cartouche `n' on line `line' of snippet `name'
\newcommand{\Cartouche}[3]{%
	{\isacartoucheopen}%
	\RawCartouche{#1}{#2}{#3}%
	{\isacartoucheclose}%
}

% \Snippet{name}: the full snippet `name'
\newcommand{\Snippet}[1]{{%
		\newcount\i
		\i=0
		\loop
		\GetSnip{#1-\the\i}%
		\advance \i 1
		\ifcsname snippet--#1-\the\i\endcsname
		\repeat
}}

% \SnippetPart{n}{m}{name}: the snippet `name' from line `n' to line `m' inclusive
\newcommand{\SnippetPart}[3]{{%
		\newcount\i
		\i=#1
		\loop
		\ifnum \i=#2
		\renewcommand{\isanewline}{}%
		\fi
		\GetSnip{#3-\the\i}%
		\advance \i 1
		\ifnum \i>#2 {}
		\else \repeat
}}


\graphicspath{images/}

% own commands

\newcommand{\IsaLogo}{\includegraphics[scale=0.07]{images/isabelle.pdf} }%

\begin{document}
	
	\title{Formalization of Thirty Three Miniatures}
	\subtitle{Notes}
	\titlerunning{Formalization of Thirty Three Miniatures}
	
	%
	% If the paper title is too long for the running head, you can set
	% an abbreviated paper title here
	%
	\author{foo}
	%
	\authorrunning{bar}
	% First names are abbreviated in the running head.
	% If there are more than two authors, 'et al.' is used.
	%
	\institute{foobar}
	%
	\maketitle              % typeset the header of the contribution      % typeset the header of the contribution
	%
	
	\section{Isabelle Questions}
	\begin{itemize}
		\item Why does Abs$\_$TYP just behave sensibly outside of its domain?
		\item What is a \enquote{qualified} lemma?
	\end{itemize}
	
	\section{Isabelle Peculiarities}
	\begin{itemize}
		\item Algebra structures are defined twice (once as type class, once as locale)
	\end{itemize}
	
	\section{Math Stuff}
	\begin{itemize}
		\item useful to realize implicit assumptions such as commutativity in a definition of sums
	\end{itemize}
	
	\section{TODO: Potentially Interesting Talk Content}
	\begin{itemize}
		\item Partially done: Overview over LA/Algebra in Isabelle (mention the chaos of everything being defined twice)
		\item Not done at all: Geometry reasoning in Isabelle
		\item Not done at all: Automation of LA proof methods
		\item Not done at all: Understand existing proof methods (argo is used a lot...?)
		\item Not done at all: ML for ?thesis in subgoals
	\end{itemize}
	
	\section{Talk Structure?}
	\begin{enumerate}
		\item Formalization content (general topic, interesting theorems)
		\item Formalization tools (Isabelle/HOL, AFP theories)
		\item Formalization struggles/observations
		\begin{itemize}
			\item Content: necessary care for detail, easier to see potential for generalization/find errors and lack of detail in proofs/...
			\item Isabelle: (Algebra) structures defined in multiple ways in different theories, reproving the same claims in the same proof twice?, excessive unfolding, no schematic variables for subgoals, ... 
		\end{itemize}
		\item Isabelle adventures:
		\begin{itemize}
			\item Eisbach proof methods
			\item Isabelle/ML?
		\end{itemize}
	\end{enumerate}
	
\end{document}